\documentclass{article}

\usepackage{amsmath,amssymb}
\usepackage{xurl}
\usepackage[hidelinks]{hyperref}
\setlength{\emergencystretch}{2em}

% Shared commands for notes in docs/paper-gaps/.

\DeclareUnicodeCharacter{2080}{\ifmmode{}_{0}\else\textsubscript{0}\fi}
\DeclareUnicodeCharacter{2081}{\ifmmode{}_{1}\else\textsubscript{1}\fi}
\DeclareUnicodeCharacter{2082}{\ifmmode{}_{2}\else\textsubscript{2}\fi}
\DeclareUnicodeCharacter{2099}{\ifmmode{}_{n}\else\textsubscript{n}\fi}

\newcommand{\Lean}{\textsc{Lean}}
\ExplSyntaxOn
\NewDocumentCommand{\leanid}{m}{
  \ifmmode
    \textsf{\detokenize{#1}}
  \else
    \group_begin:
    \urlstyle{sf}
    \tl_set:Nn \l_tmpa_tl {#1}
    \tl_replace_all:Nnn \l_tmpa_tl {\_} {_}
    \exp_args:NV \path \l_tmpa_tl
    \group_end:
  \fi
}
\ExplSyntaxOff
\providecommand{\tr}{\operatorname{tr}}
\newcommand{\tnleanIssueBase}{https://github.com/LionSR/TNLean/issues/}
\newcommand{\tnleanPrBase}{https://github.com/LionSR/TNLean/pull/}
\newcommand{\tnleanDocBase}{https://sirui-lu.com/TNLean/docs/}
\newcommand{\ghissue}[1]{\href{\tnleanIssueBase#1}{GitHub issue \##1}}
\newcommand{\ghpr}[1]{\href{\tnleanPrBase#1}{PR \##1}}
\newcommand{\doclink}[1]{\href{\tnleanDocBase#1.html}{\path{#1}}}

% Dirac notation and the identity operator, as in blueprint/src/macros/common.tex.
% \providecommand keeps note-local definitions authoritative.
\usepackage{dsfont}
\providecommand{\ket}[1]{|#1\rangle}
\providecommand{\bra}[1]{\langle #1|}
\providecommand{\braket}[2]{\langle #1 | #2 \rangle}
\providecommand{\ketbra}[2]{|#1\rangle\!\langle #2|}
\providecommand{\Id}{\mathds{1}}
\providecommand{\one}{\mathds{1}}
\providecommand{\C}{\mathbb{C}}
\providecommand{\MN}[1]{M_{#1}(\C)}
\providecommand{\MD}{M_D(\C)}

% External referencing for paper sources.  The build helper writes sanitized
% auxiliary files under build/paper-aux/ and excludes duplicated source labels.
\usepackage{xr}

% Import a generated paper auxiliary file with a caller-supplied label prefix.
% The candidate paths support builds from docs/paper-gaps/, the repository root,
% and build/paper-aux/ itself.
\newcommand{\importpaperaux}[2]{%
  \IfFileExists{../../build/paper-aux/#2.aux}{%
    \externaldocument[#1]{../../build/paper-aux/#2}%
  }{%
    \IfFileExists{build/paper-aux/#2.aux}{%
      \externaldocument[#1]{build/paper-aux/#2}%
    }{%
      \IfFileExists{#2.aux}{%
        \externaldocument[#1]{#2}%
      }{}%
    }%
  }%
}

% General external reference macros
\newcommand{\paperref}[2]{\ref{#1-#2}}
\newcommand{\papereqref}[2]{(\ref{#1-#2})}

% CPSV16 source (1606.00608 / MPDO-22-12-17-2.tex) external document and shortcuts
\importpaperaux{cpsv16-}{1606.00608/MPDO-22-12-17-2-tnlean}
\newcommand{\cpsvref}[1]{\paperref{cpsv16}{#1}}
\newcommand{\cpsveqref}[1]{\papereqref{cpsv16}{#1}}

% Verdict marker: \gapnote{<kind>}{<status>}. Kind is the policy
% classification (clarification, local-correction, scope-restriction,
% unfaithful, false-source, open-gap); status is open, wip (elimination
% underway), resolved, or historical. Machine-read by the site index;
% prints a verdict line.
\newcommand{\gapnote}[2]{\par\noindent\textbf{Verdict:} #1 (#2).\par}


\title{The Arbitrary-Chain Topological Projector Recursion}
\author{}
\date{2026-07-23}

\begin{document}
\maketitle
\gapnote{open-gap}{resolved}

\section{The source statement}

After Theorem~4.14, Cirac--P\'erez-Garc\'ia--Schuch--Verstraete apply the
fusion maps successively along an arbitrary chain. In the notation of
lines~995--1010 of Ref.~\cite{Cirac2016MPDO_arXiv}, their construction gives
\begin{align}
    \rho^{(N)}(M)
    &\simeq_{\mathrm{loc.\ iso.}}
      \mu^{\otimes N}\widetilde UQ_N\widetilde U^\dagger,
    \label{eq:cpsv16_topological_projector_recursion_source_chain_decomposition}
\end{align}
where $\simeq_{\mathrm{loc.\ iso.}}$ is the local-isometry relation of
Proposition~4.13 in Ref.~\cite{Cirac2016MPDO_arXiv}, and $\widetilde U$ is a
sequential circuit assembled from the pairwise fusion maps.

\paragraph{Local fix: rectangular vertical map.}
The retained nonzero sectors may have smaller dimension than the original
one-site space. The formal vertical map is therefore a rectangular
retained-row coisometry. Local equivalence consists of exact forward and
reconstruction identities, rather than conjugation by a square unitary. This
point is documented in
\path{docs/paper-gaps/cpgsv17_vertical_isometry_zero_sector.tex}.

The same passage of Ref.~\cite{Cirac2016MPDO_arXiv} asserts
\begin{align}
    [\mu^{\otimes N},\widetilde UQ_N\widetilde U^\dagger]
    &=0.
    \label{eq:cpsv16_topological_projector_recursion_source_commutator}
\end{align}
If the coefficients are independent of the positive chain length, every
diagonal entry of each $\chi_{\alpha,\beta,\gamma}$ lies in $\{0,1\}$.
Spectral decomposition of the terminal matrices then gives the
length-independent form stated at lines~1013--1016 of
Ref.~\cite{Cirac2016MPDO_arXiv}:
\begin{align}
    \rho^{(N)}(M)
    &=\sum_i\lambda_iP_i^{(N)}e^{-H_N},
    \label{eq:cpsv16_topological_projector_recursion_source_gibbs_decomposition}\\
    [P_i^{(N)},e^{-H_N}]
    &=0.
    \label{eq:cpsv16_topological_projector_recursion_source_gibbs_commutator}
\end{align}
Here $H_N$ is a translationally invariant commuting nearest-neighbor
Hamiltonian, and the $P_i^{(N)}$ are orthogonal projections.

\section{The formalized recursion and all-label density decomposition}

Let $P_\gamma$ be terminal orthogonal projections. One pairwise fusion step is
the direct sum
\begin{align}
    Q_{\alpha,\beta}
    &=\bigoplus_\gamma
      \chi_{\alpha,\beta,\gamma}\otimes P_\gamma.
    \label{eq:cpsv16_topological_projector_recursion_one_layer_q}
\end{align}
Each $P_\gamma$ acts on the bond space of the vertically read tensor
$M_\gamma$. The terminal matrix $\tr(M_\gamma)$ in
Ref.~\cite{Cirac2016MPDO_arXiv} closes the horizontal operator leg and leaves
these two bond indices. In the formalization, this matrix is the
physical-trace transfer of $M_\gamma$. Length independence and positivity
imply that $\chi_{\alpha,\beta,\gamma}=I$ on each retained multiplicity space,
so $Q_{\alpha,\beta}$ is an orthogonal projection. This result is formalized
for \leanid{MPOTensor.BNTFusionCoisometryFamily} and for the tensor-attached
clause \leanid{MPOTensor.BNTFusionTensorClause}.

The pairwise fusion identity is the well-typed equality
\begin{align}
    U_{\alpha,\beta}\tr(M_\alpha M_\beta)U_{\alpha,\beta}^\dagger
    &=\bigoplus_\gamma
      \chi_{\alpha,\beta,\gamma}\otimes\tr(M_\gamma).
    \label{eq:cpsv16_topological_projector_recursion_one_layer_terminal_fusion}
\end{align}
Moreover, the coisometry identity
$U_{\alpha,\beta}U_{\alpha,\beta}^\dagger=I$ implies that
$U_{\alpha,\beta}^\dagger Q_{\alpha,\beta}U_{\alpha,\beta}$ is an orthogonal
projection whenever $Q_{\alpha,\beta}$ is one. The older family
\leanid{MPOTensor.BNTFusionIsometryFamily} proves the same one-layer statement
under the additional full-support hypothesis. It is a specialization, not a
construction on a different terminal space.

\paragraph{Local fix: terminal trace orientation.}
The active fusion map has the retained-row orientation
$U_{\alpha,\beta}U_{\alpha,\beta}^\dagger=I$, as recorded in
\path{docs/paper-gaps/cpsv16_figure11_fusion_coisometry.tex}. Pull request
\#4640 temporarily interpreted the terminal trace as the one-site closed
operator $O_1(M_\gamma)$. That operator closes the bond leg and acts on the
common horizontal space. By contrast, the terminal factor in the diagram
defining $Q_N$ closes the horizontal operator leg and acts on the bond space
of $M_\gamma$. The two contractions appear separately in
\path{Papers/1606.00608/MPDO_O_L_A_.png} and
\path{Papers/1606.00608/MPDO_Structure.png}.
Equation~\ref{eq:cpsv16_topological_projector_recursion_one_layer_terminal_fusion}
restores the latter contraction.

For a matrix product density operator, every terminal matrix with this
orientation is positive semidefinite:
\begin{align}
    \tr(M_\gamma)
    &\geq 0.
    \label{eq:cpsv16_topological_projector_recursion_terminal_positivity}
\end{align}
At chain length one, the retained vertical congruence preserves positivity.
Restricting to one label $\gamma$ and one diagonal coordinate $q$ of its
positive multiplicity matrix gives the principal-block inequality
\begin{align}
    (\mu_\gamma)_{q,q}\tr(M_\gamma)
    &\geq 0.
    \label{eq:cpsv16_topological_projector_recursion_weighted_terminal_positivity}
\end{align}
Since $(\mu_\gamma)_{q,q}>0$, division by this scalar proves
Eq.~\ref{eq:cpsv16_topological_projector_recursion_terminal_positivity}.
This argument is formalized by
\leanid{MPOTensor.BNTFusionTensorClause.physTraceTransfer_verticalBNTMPO_posSemidef}.
It supplies the positivity required for the spectral decomposition at
lines~1010--1016 of Ref.~\cite{Cirac2016MPDO_arXiv}, but it does not assert
that the terminal matrices are projections.

The pairwise construction has been iterated along an arbitrary finite chain
with a fixed initial label. For total chain length $N\geq1$, the
appended-label list has length $N-1$; below, $Q_N$ and $W_N$ are indexed by
this total length. A fusion history $h$ records the intermediate labels and
multiplicity indices. Its weight $w(h)$ is the product of the selected
diagonal entries of the structure matrices, and its terminal space is the
bond space of its final label $\gamma(h)$. The resulting recursively typed
operator is
\begin{align}
    Q_N(P)
    &=\bigoplus_h w(h)P_{\gamma(h)}.
    \label{eq:cpsv16_topological_projector_recursion_recursive_q}
\end{align}

The sequential product of active fusion coisometries is formalized together
with its coisometry identity. It maps every letter of the left-associated
product tensor to
Eq.~\ref{eq:cpsv16_topological_projector_recursion_recursive_q}, with terminal
block $P_\gamma=M_\gamma^{i,j}$. Summing over equal physical indices proves
the fixed-label terminal-trace identity
\begin{align}
    W_N\tr(M_{\alpha,\boldsymbol\beta})W_N^\dagger
    &=Q_N((\tr(M_\gamma))_\gamma).
    \label{eq:cpsv16_topological_projector_recursion_fixed_label_trace}
\end{align}
Under length independence, every history weight is one. Consequently,
$Q_N(P)$ is an orthogonal projection whenever every $P_\gamma$ is an
orthogonal projection. The source embedding $\widetilde U_N$ is
$W_N^\dagger$, the adjoint of this retained-row coisometry. Exact reverse
fusion uses the active-support reconstruction in Appendix~C.4,
lines~2020--2029, of Ref.~\cite{Cirac2016MPDO_arXiv}, as documented in
\path{docs/paper-gaps/cpsv16_figure11_fusion_coisometry.tex}.

The vertical canonical form has also been assembled over every initial label
and multiplicity coordinate. For a sitewise configuration
$c=((\alpha_n,q_n))_{n=0}^{N-1}$, define
\begin{align}
    m(c)
    &=\prod_{n=0}^{N-1}(\mu_{\alpha_n})_{q_n,q_n}.
    \label{eq:cpsv16_topological_projector_recursion_chain_weight}
\end{align}
The first label is the initial label in the recursion, and the remaining
$N-1$ labels form the appended list. If $W_{N,c}$ is the resulting
retained-row sequential coisometry, define
\begin{align}
    Q_{N,c}
    &=Q_N((\tr(M_\gamma))_\gamma).
    \label{eq:cpsv16_topological_projector_recursion_configuration_terminal_q}
\end{align}
The retained density block is
\begin{align}
    D_{N,c}
    &=m(c)W_{N,c}^\dagger Q_{N,c}W_{N,c}.
    \label{eq:cpsv16_topological_projector_recursion_retained_density_block}
\end{align}
Thus the embedding $\widetilde U$ in
Eq.~\ref{eq:cpsv16_topological_projector_recursion_source_chain_decomposition}
is the adjoint $W_{N,c}^\dagger$ of the coisometry used in the formalization.
This orientation agrees with
\path{docs/paper-gaps/cpsv16_figure11_fusion_coisometry.tex}.

Let $V$ be the rectangular vertical coisometry of Proposition~4.13 in
Ref.~\cite{Cirac2016MPDO_arXiv}, and let $R_N$ be the direct sum of the blocks
in Eq.~\ref{eq:cpsv16_topological_projector_recursion_retained_density_block}
over all configurations $c$. For every $N>0$, the formalization proves
\begin{align}
    V^{\otimes N}\rho^{(N)}(M)(V^{\otimes N})^\dagger
    &=R_N,
    \label{eq:cpsv16_topological_projector_recursion_all_label_density_forward}\\
    ((V^\dagger)^{\otimes N})R_NV^{\otimes N}
    &=\rho^{(N)}(M).
    \label{eq:cpsv16_topological_projector_recursion_all_label_density_reverse}
\end{align}

\paragraph{Local fix: rectangular vertical map.}
The reverse identity
Eq.~\ref{eq:cpsv16_topological_projector_recursion_all_label_density_reverse}
uses the exact one-site reconstruction from the retained nonzero sectors in
Proposition~4.13 and Appendix~C.4, lines~2020--2029, of
Ref.~\cite{Cirac2016MPDO_arXiv}. It does not require full support or a square
unitary. The rectangular zero-sector issue is documented in
\path{docs/paper-gaps/cpgsv17_vertical_isometry_zero_sector.tex}.
Equation~\ref{eq:cpsv16_topological_projector_recursion_chain_weight} is
exactly the selected diagonal coefficient of $\mu^{\otimes N}$. Hence
Eqs.~\ref{eq:cpsv16_topological_projector_recursion_all_label_density_forward}
and~\ref{eq:cpsv16_topological_projector_recursion_all_label_density_reverse}
prove the density identity at line~999 of
Ref.~\cite{Cirac2016MPDO_arXiv}, with source chain length $N$ and no additional
compatibility hypothesis.

The assembled operator is
\leanid{MPOTensor.BNTFusionTensorClause.topologicalDensityOperatorSucc}. Its
argument is one less than the positive source chain length. The source-facing
existence theorem
\leanid{MPOTensor.exists_topologicalDensityDecomposition_of_isRFPViaTS}
chooses one vertical decomposition and proves both all-label identities for
every positive $N$.

\section{The all-label commutator}

In the same all-label coordinates, define
\begin{align}
    A_N
    &=\bigoplus_c m(c)I_c,
    \label{eq:cpsv16_topological_projector_recursion_multiplicity_factor}\\
    T_N
    &=\bigoplus_c W_{N,c}^\dagger Q_{N,c}W_{N,c}.
    \label{eq:cpsv16_topological_projector_recursion_recursive_factor}
\end{align}
The retained operator is $R_N=A_NT_N$. Within each copy-configuration block,
$A_N$ is a scalar multiple of the identity. Direct-sum multiplication
therefore gives
\begin{align}
    A_NT_N
    &=T_NA_N.
    \label{eq:cpsv16_topological_projector_recursion_density_commutator}
\end{align}
This proves the source commutator in
Eq.~\ref{eq:cpsv16_topological_projector_recursion_source_commutator} for
every positive chain length, without length independence or a
caller-supplied compatibility condition. The two factors are
\leanid{MPOTensor.BNTFusionTensorClause.topologicalMultiplicityWeightFactorSucc}
and
\leanid{MPOTensor.BNTFusionTensorClause.topologicalRecursiveFactorSucc}.
The source-facing theorem
\leanid{MPOTensor.exists_topologicalDensityDecomposition_and_factorCommutator_of_isRFPViaTS}
chooses the same fusion clause for the all-label density identities and
Eq.~\ref{eq:cpsv16_topological_projector_recursion_density_commutator}.

\section{The terminal spectral refinement}

Fix the BNT fusion clause selected from the horizontal-canonical MPDO and its
renormalization-fixed-point data. The length-independence hypothesis concerns
the positive trace-power coefficients of this fixed clause; it does not
permit an existential choice of a different fusion clause.

For each terminal label $\gamma$, define $K_\gamma=\tr(M_\gamma)$. Since
$K_\gamma$ is positive semidefinite, it has a spectral decomposition
\begin{align}
    K_\gamma
    &=\sum_k\lambda_{\gamma,k}E_{\gamma,k},
    \label{eq:cpsv16_topological_projector_recursion_terminal_spectral_decomposition}\\
    \lambda_{\gamma,k}
    &\geq0.
    \label{eq:cpsv16_topological_projector_recursion_terminal_eigenvalue_positivity}
\end{align}
One-site MPDO positivity, restricted to a fixed BNT copy and divided by its
strictly positive multiplicity weight, proves $K_\gamma\geq0$. The operators
$E_{\gamma,k}$ are rank-one orthogonal projections.

Extend $E_{\gamma,k}$ by zero on all other terminal labels, and transport it
through the recursive fusion histories and the adjoint sequential
coisometries. The all-label assembly gives operators
$P_{\gamma,k}^{(N)}$ satisfying
\begin{align}
    (P_{\gamma,k}^{(N)})^\dagger
    &=P_{\gamma,k}^{(N)},
    \label{eq:cpsv16_topological_projector_recursion_refined_projector_self_adjoint}\\
    (P_{\gamma,k}^{(N)})^2
    &=P_{\gamma,k}^{(N)},
    \label{eq:cpsv16_topological_projector_recursion_refined_projector_idempotent}\\
    P_{\gamma,k}^{(N)}P_{\delta,\ell}^{(N)}
    &=0
    \quad\text{if }(\gamma,k)\neq(\delta,\ell).
    \label{eq:cpsv16_topological_projector_recursion_refined_projector_orthogonality}
\end{align}
For every positive chain length,
\begin{align}
    T_N
    &=\sum_{\gamma,k}\lambda_{\gamma,k}P_{\gamma,k}^{(N)},
    \label{eq:cpsv16_topological_projector_recursion_terminal_spectral_refinement}\\
    [A_N,P_{\gamma,k}^{(N)}]
    &=0,
    \label{eq:cpsv16_topological_projector_recursion_projector_weight_commutator}\\
    R_N
    &=\sum_{\gamma,k}
      \lambda_{\gamma,k}A_NP_{\gamma,k}^{(N)}.
    \label{eq:cpsv16_topological_projector_recursion_terminal_density_refinement}
\end{align}
This refinement is formalized by
\leanid{MPOTensor.terminalSpectralProjectorRefinement_of_isRFPViaTS}.

\section{Remaining physical-space construction}

The all-label density identity corresponding to
Eq.~\ref{eq:cpsv16_topological_projector_recursion_source_chain_decomposition},
the commutator in
Eq.~\ref{eq:cpsv16_topological_projector_recursion_source_commutator}, and the
terminal spectral refinement in
Eq.~\ref{eq:cpsv16_topological_projector_recursion_terminal_spectral_refinement}
are combined with an explicit commuting nearest-neighbor Hamiltonian on the
retained vertical space. For every chain length $N\geq2$, the two-site
interaction is the logarithmic multiplicity energy on its first site tensored
with the identity on its second site. Its periodic translates are diagonal
and commute, and their negative exponential is $\mu^{\otimes N}$. The
terminal spectral family has cardinality at most the physical one-site
dimension $d$ and is extended by zero to exactly $d$ indices. This proves the
retained-coordinate analogue of
Eqs.~\ref{eq:cpsv16_topological_projector_recursion_source_gibbs_decomposition}
and~\ref{eq:cpsv16_topological_projector_recursion_source_gibbs_commutator}.

The adjoint retained-row map reconstructs the physical density operator from
this analogue. It does not yet produce physical projectors and a finite
physical-space Hamiltonian because the rectangular vertical map may have a
nontrivial complementary zero sector. The required complementary-subspace
extension is recorded in
\path{docs/paper-gaps/cpsv16_topological_gibbs_physical_complement.tex}.

Independently, Ref.~\cite{Cirac2016MPDO_arXiv} defines the local term on two
spins but gives no separate convention for a one-site periodic chain. This
boundary is recorded in
\path{docs/paper-gaps/cpsv16_topological_gibbs_length_one.tex}.

\bibliographystyle{alpha}
\bibliography{references}

\end{document}
